\chapter{Conclusion}

This thesis focused on the search for ultra high energy neutrinos in ice via Askaryan emission.
These neutrinos originate in extreme astrophysical sources such as active galactic nuclei or, cosmogenically, when ultra high energy cosmic rays interact with the cosmic microwave background. 
My doctoral work focused on simulating present and potential neutrino detectors; developing, assembling, and maintaining detector hardware; and searching for neutrinos in over one decade of experimental data. 

This thesis began by presenting a study simulating radio antennas accompanying all 9,760 optical modules in the proposed IceCube Gen2 optical volume. 
These were intended to enhance IceCube’s directional reconstruction of neutrino cascades with Askaryan emission - cascades from neutrinos with more than 6~PeV of energy. 
These events are very rare, therefore maximizing the information captured by their observation is crucial. 
Due to the thin width of the Askaryan emission cone and low signal strength, only 11\% of the 54 expected IceCube Gen2 Glashow Resonance events had Askaryan content in 10 years of observations. 
This fraction increases for higher energy neutrinos, but the expected number of higher energy neutrinos falls too quickly to see much benefit.
Given that most observed events will have low radio content, it may be more lucrative to use strategically placed phased arrays in the IceCube Gen2 optical module, though a simulation study is necessary to assess potential benefits of this. 
This study was completed with Lu Lu, Albrecht Karle, and Ben Hokanson-Fasig \citep{icecube_gen2_infill}. 

This thesis then moved into discussions of the  Askaryan Radio Array (ARA). 
ARA is comprised of 5 radio stations and has taken data since 2011. 
Updates to simulations of the array were discussed, highlighting the importance of simulating the outgoing muons and taus from charge current neutrino interactions of the same flavors and the benefits of simulating these events in an array-wide context that considers an events ability to trigger multiple stations. 
We determined that, at trigger level, ARA is the most sensitive detector for neutrinos with energy above 3 EeV and could have observed 2.27 neutrinos according to the most optimistic \cite{flux-kotera} neutrino model. 
These trigger level figure demonstrate the value of the ARA dataset and the goal that neutrino searches should strive for. 
These studies were conducted with Austin Cummings, Ryan Krebs, Marco Muzio, and Alan Salcedo Gomez \citep{ara_secondaries, ara_veff}. 

The thesis finishes with a neutrino search through 10.6 years of ARA data. 
This was performed by a team of 8 - myself, Brian Clark, Paramita Dasgupta, Pawan Giri, Mohammad Hossain, Alex Machtay, Alan Salcedo Gomez, and was led by Marco Muzio \citep{ara_5sa-icrc}. 
My work focused on the cleaning of ARA's second station and simulations for the analysis. 
The analysis had one passing event on the PA station that was later determined to be a glitch which has a dedicated analysis cut for the traditional stations A1-A5 but not for the PA detector as it has a different kind of data acquisition board from the traditional stations. 
We've since determined the glitch does not occur in the data acquisition boards, but is a power surge near the in-ice antennas causing amplified waveforms in one string of antennas and releasing radio frequency emission that is recorded by nearby strings. 
As A5 is colocated with PA, the passing PA event was the observation of a glitch on A5. 
Our analysis set the upper limit on neutrino sensitivity above 40 EeV and was expected to observe 0.41 events according to the \cite{flux-kotera} model, 18\% less than the estimate at trigger level. 
This is mainly due to the inefficient cuts that remove surface related events, mainly human-made anthropogenic backgrounds that mimic neutrinos and reconstruct deep into the ice. 
Future analyses should consider ways to improve the efficiency of cuts targeting surface backgrounds, perhaps considering different cuts for the busy Pole season and the quieter winterover period. % I think we did this and determined we shouldn't? Why?
We also recommend installing detectors more than 5 kilometers from expected anthropogenic backgrounds. 
This analysis also expected 1.75 events from the \cite{km3net_neutrino} flux and found 0. 
Finally, the analysis involved a great number of station-level and array-wide configurations which drove up simulation and analysis optimization complexity. 
We recommend that present and future detectors minimize changes that will motivate the creation of new station configurations such as updates to station hardware and data taking settings. 
The gold standard would be installing a detector in an optimized state and leaving it as long as possible as such. 
If a change is necessary, survey all stations for necessary changes and update them all at once, so the array wide configuration is only updated once. 

Thesis work that we're not presented in the main body of this thesis include the design of antenna centralizers that are deployed in the Radio Neutrino Observatory of Greenland stations, the GPS survey of ARA, and my work configuring the ARA simulation software to receive electric field at each antenna as input which allows the integration of softwares like Corsika and FAERIE with ARA detector simulations. 